\chapter*{Conclusion et perspectives}
\addcontentsline{toc}{chapter}{Conclusion et perspectives}
\markboth{Conclusion}{Conclusion}
\label{sec:conclusion}

Le projet \textbf{Tarteel Maghreb} a consisté à concevoir et développer une application mobile innovante destinée à l’apprentissage de la récitation coranique selon la lecture \textbf{Qālūn ‘an Nāfi‘}. Ce travail nous a permis de mettre en pratique les compétences acquises en développement mobile, en développement back-end ainsi qu’en intégration de solutions d’intelligence artificielle. \vspace{2mm}

La démarche suivie a été structurée et progressive. Elle a débuté par une \textbf{analyse des besoins}, comprenant l’étude de l’existant, l’identification des acteurs et la définition des exigences fonctionnelles et non fonctionnelles. Elle s’est poursuivie par une phase de \textbf{conception}, avec l’élaboration d’une architecture en trois tiers, la modélisation des vues statiques et dynamiques, et la préparation d’interfaces intuitives pour l’utilisateur. Enfin, la phase de \textbf{réalisation} a permis d’implémenter les fonctionnalités principales, d’intégrer la reconnaissance vocale via le modèle \textit{Whisper}, de développer l’API back-end avec \textit{Node.js} et un stockage cloud sécurisé, et de créer une application mobile ergonomique et fluide. 
\vspace{2mm}

Le résultat final est une application fonctionnelle et accessible, capable de fournir un retour pédagogique intelligent sur la récitation et d’accompagner l’utilisateur dans sa progression. Cette réalisation illustre la synergie entre tradition et technologie et constitue une base solide pour le développement futur d’outils pédagogiques coraniques.\vspace{2mm}

Plusieurs \textbf{perspectives d’évolution} sont envisageables pour enrichir l’application. Il s’agit notamment de l’intégration d’autres riwāyah pour élargir le public, de l’optimisation du module d’intelligence artificielle pour améliorer la précision de la détection des erreurs de tajwīd et proposer un retour plus détaillé, de l’introduction de mécanismes de gamification (défis, badges, classements) pour accroître l’engagement des utilisateurs, ainsi que du renforcement du support multilingue et de l’accessibilité. D’autres pistes incluent la synchronisation et la sauvegarde cloud pour assurer la continuité de l’apprentissage sur plusieurs appareils, ainsi que l’exploitation avancée des données afin de proposer des parcours personnalisés et adaptatifs.\vspace{2mm} 

Ainsi, \textbf{Tarteel Maghreb} constitue une contribution significative au domaine des technologies éducatives appliquées au Coran et ouvre la voie à de nombreuses perspectives pour améliorer l’expérience des apprenants et renforcer l’impact pédagogique de l’application.
